IoT devices send tasks with different latency limits to a small pool of edge servers. Deep reinforcement learning (RL) is often used to decide where each task runs, and many methods give every missed task the same penalty. This paper studies the per-class penalty against an exact optimum. Four offloading systems with up to four heterogeneous edge nodes, three task classes, random background traffic and a cloud that meets a limit only by chance are solved by finite-horizon dynamic programming (DP). DP gives the optimal policy, the optimal trade-off between the classes and an exact score for any trained agent, so no result depends on a test set. A short proof shows that a reward with a bonus for a met limit ranks policies by the penalty plus the bonus. A class-blind penalty therefore minimizes the overall miss rate, and its optimal policy is 65\% to 118\% above the optimal SLA cost. A class-weighted penalty is 1\% to 10\% above it, and a corrected penalty that folds in the bonus reaches it. PPO, DQN and A2C agents get the same tuning budget and are trained with 10 seeds. In the nominal system, PPO with the weighted reward misses 2.7\% of critical tasks against 11.7\% with the class-blind reward, is 10\% above the optimum, and beats a reservation rule tuned on the exact model in 9 of 10 seeds. In three larger systems the weighted reward lowers the SLA cost by 18\% to 63\% against the class-blind reward, but the agents stay 36\% to 67\% above the optimum, and a rule with per-class node sets tuned on the exact model is better. A rebuilt published design with adaptive penalty tiers is 41\% to 60\% above the optimum, because its penalties drift together when its targets cannot all be met. All code and results are open.
